
\documentclass{article}
\usepackage{colortbl}
\usepackage{makecell}
\usepackage{multirow}
\usepackage{supertabular}

\begin{document}

\newcounter{utterance}

\centering \large Interaction Transcript for game `referencegame', experiment `line\_grids\_columns', episode 3 with reflexion{-}Qwen3.5{-}4B.
\vspace{24pt}

{ \footnotesize  \setcounter{utterance}{1}
\setlength{\tabcolsep}{0pt}
\begin{supertabular}{c@{$\;$}|p{.15\linewidth}@{}p{.15\linewidth}p{.15\linewidth}p{.15\linewidth}p{.15\linewidth}p{.15\linewidth}}
    \# & $\;$A & \multicolumn{4}{c}{Game Master} & $\;\:$B\\
    \hline

    \theutterance \stepcounter{utterance}  
    & & \multicolumn{4}{p{0.6\linewidth}}{
        \cellcolor[rgb]{0.9,0.9,0.9}{
            \makecell[{{p{\linewidth}}}]{
                \texttt{\tiny{[P1$\langle$GM]}}
                \texttt{You have attempted this same task before and did not succeed. These are the plans you wrote after those attempts.} \\
\texttt{Trial \#1: New plan: My previous description was too vague because it failed to distinguish the target grid from the distractors. Specifically, I did not mention the specific pattern of the empty cells in the middle columns, which is the key difference between the target (a solid block of 3 empty columns) and Distractor 2 (a scattered pattern of empty cells). Next time, I will describe the target grid by explicitly stating that the first and fifth columns are filled with 'X' and the middle three columns are completely empty, ensuring this unique structural feature is highlighted to separate it from the other grids.} \\
\\ 
\texttt{Trial \#2: New plan: My previous description failed because it did not explicitly quantify the number of empty columns, which is the definitive feature distinguishing the target grid (3 empty columns) from Distractor 2 (2 empty columns). Next time, I will describe the target grid by specifying that it has exactly three consecutive empty columns in the center, flanked by two filled columns, ensuring the count of empty columns is clearly stated to avoid confusion with the distractors.} \\
\\ 
\\ 
\texttt{You are given three grids, where each of them is 5 by 5 in size.} \\
\texttt{Grids have empty cells marked with "▢" and filled cells marked with "X".} \\
\texttt{Your task is to generate a referring expression that best describes the target grid while distinguishing it from the two other distractor grids.} \\
\texttt{The first grid is the target grid, and the following two grids are the distractors.} \\
\\ 
\texttt{Target grid:} \\
\\ 
\texttt{X □ □ X X} \\
\texttt{X □ □ X X} \\
\texttt{X □ □ X X} \\
\texttt{X □ □ X X} \\
\texttt{X □ □ X X} \\
\\ 
\texttt{Distractor grid 1:} \\
\\ 
\texttt{□ □ X X X} \\
\texttt{□ □ X X X} \\
\texttt{□ □ X X X} \\
\texttt{□ □ X X X} \\
\texttt{□ □ X X X} \\
\\ 
\texttt{Distractor grid 2:} \\
\\ 
\texttt{X □ X X □} \\
\texttt{X □ X X □} \\
\texttt{X □ X X □} \\
\texttt{X □ X X □} \\
\texttt{X □ X X □} \\
\\ 
\texttt{Instruction: Describe the target grid.} \\
\texttt{Generate the referring expression starting with the tag "Expression: " for the given target grid. Omit any other text.} \\
            }
        }
    }
    & \\ \\

    \theutterance \stepcounter{utterance}  
    & \multicolumn{4}{p{0.6\linewidth}}{
        \cellcolor[rgb]{0.9,0.9,0.9}{
            \makecell[{{p{\linewidth}}}]{
                \texttt{\tiny{[P1$\rangle$GM]}}
                \texttt{Expression: The grid with exactly three consecutive empty columns in the center, flanked by two filled columns on the left and right.} \\
            }
        }
    }
    & & \\ \\

    \theutterance \stepcounter{utterance}  
    & & & \multicolumn{2}{p{0.3\linewidth}}{
        \cellcolor[rgb]{0.9,0.9,0.9}{
            \makecell[{{p{\linewidth}}}]{
                \texttt{\tiny{[GM$|$GM]}}
                \texttt{Expression: The grid with exactly three consecutive empty columns in the center, flanked by two filled columns on the left and right.} \\
            }
        }
    }
    & & \\ \\

    \theutterance \stepcounter{utterance}  
    & & \multicolumn{4}{p{0.6\linewidth}}{
        \cellcolor[rgb]{0.9,0.9,0.9}{
            \makecell[{{p{\linewidth}}}]{
                \texttt{\tiny{[GM$\rangle$P2]}}
                \texttt{You have attempted this same task before and did not succeed. These are the plans you wrote after those attempts.} \\
\texttt{Trial \#1: New plan: I misinterpreted the expression "middle three columns are empty" as referring to the second, third, and fourth columns, which matches the "Second" grid. However, the expression actually describes the first and fifth columns being filled while the middle three are empty, which corresponds to the "Third" grid. Next time, I will carefully verify that the filled columns are exactly the first and fifth, and the empty columns are the middle three, before selecting the grid.} \\
\\ 
\texttt{Trial \#2: New plan: I misidentified the pattern in the "Third" grid, assuming the single filled columns on the sides represented a "block" or that the description allowed for a different configuration. The expression specifically requires a "solid block of three empty columns in the middle," which implies the first and last columns must be filled. While the "Third" grid has filled columns on the sides, the "Second" grid has a solid block of three filled columns in the middle, which is the opposite of the description. I need to re{-}read the expression carefully: "solid block of three empty columns in the middle." This means columns 2, 3, and 4 must be empty (▢), and columns 1 and 5 must be filled (X). Looking at the grids again: Grid 1 has filled columns 1 and 5, but columns 2, 3, 4 are empty. Grid 2 has empty columns 1 and 5, but columns 2, 3, 4 are filled. Grid 3 has filled columns 1 and 5, but columns 2, 4 are empty and 3 is filled. Therefore, the correct grid is the first one. My previous mistake was selecting the third grid, likely due to a misreading of the grid layout or the expression's requirement for a "solid" block of emptiness. Next time, I will explicitly check if columns 2, 3, and 4 are all empty and columns 1 and 5 are all filled before answering.} \\
\\ 
\texttt{Answer: 1} \\
\\ 
\\ 
\texttt{You are given three grids, where each of them is 5 by 5 in size.} \\
\texttt{Grids have empty cells marked with "▢" and filled cells marked with "X".} \\
\texttt{You are also given a referring expression that describes one of the given grids.} \\
\texttt{Your task is to select the grid that matches the given referring expression.} \\
\texttt{Generate only the number (in text) of the grid that matches the given expression by selecting first, second, or third.} \\
\\ 
\texttt{First:} \\
\\ 
\texttt{X □ □ X X} \\
\texttt{X □ □ X X} \\
\texttt{X □ □ X X} \\
\texttt{X □ □ X X} \\
\texttt{X □ □ X X} \\
\\ 
\texttt{Second:} \\
\\ 
\texttt{□ □ X X X} \\
\texttt{□ □ X X X} \\
\texttt{□ □ X X X} \\
\texttt{□ □ X X X} \\
\texttt{□ □ X X X} \\
\\ 
\texttt{Third:} \\
\\ 
\texttt{X □ X X □} \\
\texttt{X □ X X □} \\
\texttt{X □ X X □} \\
\texttt{X □ X X □} \\
\texttt{X □ X X □} \\
\\ 
\texttt{Expression: The grid with exactly three consecutive empty columns in the center, flanked by two filled columns on the left and right.} \\
\texttt{Question: Which grid does the expression refer to?} \\
\texttt{Start with the tag "Answer: ", followed by your selection. Omit any other text.} \\
            }
        }
    }
    & \\ \\

    \theutterance \stepcounter{utterance}  
    & & & \multicolumn{4}{p{0.6\linewidth}}{
        \cellcolor[rgb]{0.9,0.9,0.9}{
            \makecell[{{p{\linewidth}}}]{
                \texttt{\tiny{[GM$\langle$P2]}}
                \texttt{Answer: 1} \\
            }
        }
    }
     \\ \\

    \theutterance \stepcounter{utterance}  
    & & & \multicolumn{2}{p{0.3\linewidth}}{
        \cellcolor[rgb]{0.9,0.9,0.9}{
            \makecell[{{p{\linewidth}}}]{
                \texttt{\tiny{[GM$|$GM]}}
                \texttt{1} \\
            }
        }
    }
    & & \\ \\

\end{supertabular}
}

\end{document}
